% TLDR: State and derive the proper ordinal objective and its relation to distribution error.
\section{A Distribution-Sensitive Objective}
\label{sec:objective}
Proper scoring rules reward accurate probability forecasts in expectation \citep{gneiting2007strictly}.
For the cumulative probabilities $Q_c=\sum_{r=0}^c q_r$, we use the loss
\begin{equation}
 S(q,r)=\frac{1}{5}\sum_{c=0}^{4}
       \bigl(Q_c-\mathbf{1}\{r\leq c\}\bigr)^2.
 \label{eq:rps}
\end{equation}
The factor $1/5$ places RPS on $[0,1]$.
The ordered thresholds retain the ordinal structure: probability placed many categories away from an observed rating incurs more cumulative error than probability placed in an adjacent category.
Equal threshold weights are an explicit design choice, not a claim that every practical consequence of a rating is equally spaced.

Let $P_c=\Pr(R\leq c\mid z)$.
For each threshold, expansion of the squared error gives
\begin{align}
 \mathbb E[S(q,R)\mid z]
 &=D(q,p^*)+U(p^*),\\
 D(q,p^*)&=\frac{1}{5}\sum_{c=0}^{4}(Q_c-P_c)^2,\\
 U(p^*)&=\frac{1}{5}\sum_{c=0}^{4}P_c(1-P_c).
 \label{eq:decomposition}
\end{align}
The first term measures conditional distribution error; the second is the loss remaining even for a perfect forecast of a disagreeing expert population.
Because the cumulative probabilities determine all six category probabilities, the unrestricted expected-loss minimizer is uniquely $q=p^*$.
This fact justifies the population training target.
It does not guarantee that a particular model class, prompt search, or finite sample will attain that target.

For equal category spacing on $[0,1]$, the one-dimensional Wasserstein distance is
\begin{equation}
 W_1(q,p^*)=\frac{1}{5}\sum_{c=0}^{4}|Q_c-P_c|.
\end{equation}
Cauchy--Schwarz gives $W_1(q,p^*)^2\leq D(q,p^*)$.
Thus small excess expected RPS bounds an interpretable distribution distance.
More generally, for a normalized $L$-Lipschitz utility function $f$, the expectation gap is at most $L W_1$.
This connects a distributional claim to uses such as estimating the mean rating, without requiring a separate match of every moment.

\subsection{Why correlation or a pooled divergence is insufficient}
Consider human scores $(0,0.25,0.5)$ and judge scores $(0.5,0.75,1)$.
Pearson and rank correlations are one, although every score differs by $0.5$.
Alternatively, permuting predictions among artifacts preserves the marginal score histogram while destroying itemwise agreement.
A divergence between pooled histograms misses that failure.
The required comparison is conditional on the same artifact, then averaged across artifacts.

Logarithmic loss is another proper objective and connects to Kullback--Leibler divergence, but it ignores category order and diverges when the judge assigns zero probability to an observed category.
With few ratings per artifact, directly estimating an entire human probability vector and plugging it into a divergence is also noisy.
RPS can instead score each individual rating without declaring a sparse empirical histogram to be the true human distribution.
We use cross-entropy as a separate training condition and a shared secondary evaluation metric, with the probability-floor policy specified before evaluation.

\subsection{Training objectives and artifact weighting}
For $n$ artifacts, with assigned rater set $\mathcal H_i$ of size $m_i$, define
\begin{equation}
 \widehat L_S(\theta)=\frac1n\sum_{i=1}^{n}
       \frac1{m_i}\sum_{h\in\mathcal H_i}S(q_\theta(z_i),R_{ih}).
 \label{eq:empirical}
\end{equation}
Artifacts receiving additional ratings retain the same total weight as other artifacts.
Let $\bar r_i$ be the mean normalized human rating for artifact $i$.
The correlation-training condition maximizes Pearson correlation between $\mu_q(z_i)$ and $\bar r_i$ across training artifacts.
The proper-loss-training condition minimizes Equation~\ref{eq:empirical} on the same training artifacts, with RPS primary.
The judge class, inputs, data access and search budget are matched between conditions.
Pearson correlation constrains the relationship between output ratings; the mapping from source artifacts to predictions may be nonlinear.
Constant predictions or targets make correlation undefined and require a prespecified failure policy.

Additional proper-loss conditions train with cross-entropy or six-category Brier loss and remain separately identified.
Cross-entropy scores an individual rating as $-\log q_r$ using natural logarithms.
The six-category Brier loss is $\sum_{k=0}^{5}(q_k-\mathbf1\{r=k\})^2$, with range $[0,2]$.
It differs from the binary Brier loss used for a separate prover-outcome target.
No population optimality result proves that one training procedure will perform better with finite data and a restricted judge class.

Every condition selects hyperparameters by the same validation RPS.
Changing the validation criterion from correlation to proper loss would answer a different question from changing the training objective.
All prompt changes, fitted mappings and attempted candidates are recorded because repeated selection can overfit validation data~\citep{cawley2010overfitting}.
Final test ratings do not enter training or selection.

\subsection{Shared evaluation metrics}
RPS is the primary held-out metric for every training condition.
Cross-entropy and total variation distance (TVD) are shared secondary metrics.
For an artifact's empirical human distribution $\widehat p_i$, define
\begin{equation}
 \widehat{\operatorname{TVD}}=\frac1n\sum_i\frac12
       \sum_{r=0}^{5}|q_\theta(r\mid z_i)-\widehat p_i(r)|.
\end{equation}
This TVD lies in $[0,1]$ and weights artifacts equally.
It is a finite-rating distribution diagnostic, not a proper individual-rating loss or total-variation mutual information.
The empirical human distribution remains noisy when few people rate an artifact.
Expected-rating correlations, absolute error, signed bias and human--human agreement supplement the common loss panel.

\subsection{A worked distinction between uncertainty sources}
Suppose, illustratively, that half of the qualified raters assign category zero to one artifact and half assign category five.
The correct human forecast places probability $1/2$ on each endpoint, so every cumulative probability $P_c$ equals $1/2$ for $c=0,\ldots,4$.
Its expected RPS is $0.25$, entirely due to human disagreement.
A judge that always assigns all probability to category zero has expected RPS $0.50$: its additional distribution error is $0.25$.
Both forecasts can be perfectly reproducible across repeated judge calls.
The deterministic endpoint forecast is therefore stable but worse at predicting the human distribution.

These numbers are an exact illustrative calculation, not observed study results.
They show why a low model-repeat standard deviation cannot replace the proper-loss comparison and why a positive raw RPS cannot by itself imply a defective judge.
The same distinction applies when interpreting the dense subset: additional raters estimate the human cumulative probabilities more precisely, while additional judge calls estimate the variability of the probability forecast.
The two replication budgets answer different questions and must be reported separately.
